\documentclass[11pt]{amsproc}

\usepackage[utf8]{inputenc}
\usepackage[T1]{fontenc}
\usepackage[spanish,es-noshorthands,es-nodecimaldot]{babel}
\usepackage{amsmath,amssymb}
\usepackage[margin=2.5cm]{geometry}
\usepackage{enumitem}

\begin{document}

\section*{Problema} 

Encuentre pares de números racionales
$(a_0,b_0),(a_1,b_1),\dots$, en los que $a_i$ y $b_i$ nunca son ambos nulos,
con la siguiente propiedad: para todo par racional $(\alpha,\beta)\neq(0,0)$
existe un $N$ tal que cualesquiera $N$ pares consecutivos contienen un par
$(a_j,b_j)$ con $\alpha a_j+\beta b_j=0$.

\section*{Solución} 

Para $(\alpha,\beta)\neq(0,0)$
racional, sea
\[
  \ell(\alpha,\beta)=\{(a,b):\alpha a+\beta b=0\}.
\]
Hay una cantidad numerable de pares $(\alpha,\beta)$, por lo que hay una
cantidad numerable de rectas, digamos 
$\ell_1,\ell_2,\ell_3,\dots$. Las repeticiones no causan ningún problema. 

Cada recta contiene un punto racional no nulo. Si
$\ell_n=\ell(\alpha,\beta)$, tomamos $q_n=(\beta,-\alpha)$. Entonces
$\alpha\beta+\beta(-\alpha)=0$, y $q_n\neq(0,0)$.

Todo entero $m\geq1$ se escribe
así: $m=2^v\cdot u$ con $u$ impar, $v=v_2(m)$.
Definimos
\[
  (a_i,b_i)=q_{\,v_2(i+1)+1}\qquad(i\geq0).
\]
Así, $(a_i,b_i)$ está en $\ell_n$ exactamente cuando $v_2(i+1)=n-1$. Leyendo
$n=v_2(i+1)+1$ para $i=0,1,2,\dots$ se obtiene la sucesión 
\[
  1,2,1,3,1,2,1,4,1,2,1,3,1,2,1,5,\dots
\]
Cada par es algún $q_n$, así que ningún par es $(0,0)$.

Tomemos ahora $2^n$ índices consecutivos $i,i+1,\dots,i+2^n-1$. Los números
$i+1,\dots,i+2^n$ son $2^n$ enteros consecutivos, de modo que recorren cada
clase módulo $2^n$ exactamente una vez. En particular, uno de
ellos, digamos $j+1$, satisface
\[
  j+1\equiv 2^{n-1}\bmod{2^n}, 
\]
es decir $j+1=2^{n-1}(1+2k)$ para algún $k\ge 0$, así que $v_2(j+1)=n-1$. Por
lo tanto, $(a_j,b_j)=q_n\in\ell_n$.

Ya casi estamos. Sea $(\alpha,\beta)\neq(0,0)$ racional.
Su recta $\ell(\alpha,\beta)$ es $\ell_n$ para algún $n$. Tomemos $N=2^n$.
Sabemos que $N$ pares consecutivos contienen un par con
$\alpha a_j+\beta b_j=0$. 

\end{document}
